\chapter{Unconditional Resolution Width Lower Bound: $\Omega(n)$}

\section{The Median Boundary Lemma}
\begin{lemma}[Existence of the Median Clause $C^*$]
In every Resolution refutation $\Pi = (C_1, \dots, C_S = \bot)$ of $\Phi_X$, there exists at least one clause $C^* \in \Pi$ satisfying:
\[
\frac{1}{3}\mu_0 n \le \mu(C^*) \le \frac{2}{3}\mu_0 n.
\]
\end{lemma}

\begin{proof}
At the start, all initial axioms $A \in \Phi_X$ have $\mu(A) = 1$.
At the end, the empty clause $\bot$ requires a full 2-cycle to semantically contradict, so $\mu(\bot) \ge \mathrm{Sys}_2(X) \ge \mu_0 n$.
Because resolution steps are binary, $\mu(C) \le \mu(A) + \mu(B) \le 2 \max(\mu(A), \mu(B))$.
Thus, the measure $\mu$ cannot jump over an interval of length $\ge \frac{1}{3}\mu_0 n$ in a single step without landing inside $[\frac{1}{3}\mu_0 n, \frac{2}{3}\mu_0 n]$.
\end{proof}

\section{The Resolution Width Theorem}
\begin{theorem}[Unconditional Resolution Width on Ramanujan 2-Complexes]
For any $(n, k, \gamma)$-LSV complex $X$, the refutation width satisfies:
\[
\Width(\Phi_X \vdash \bot) \ge \frac{1}{3} \gamma \mu_0 \cdot n = \Omega(n).
\]
\end{theorem}

\begin{proof}
Let $C^*$ be the median clause from Lemma 5.1.
By Chapter 4, Theorem 4.1, because $\mu(C^*) \le \frac{2}{3}\mu_0 n$, the boundary retention invariant gives:
\[
|C^*| = |\Vars(C^*)| \ge \gamma \cdot \mu(C^*) \ge \gamma \cdot \left(\frac{1}{3}\mu_0 n\right) = \frac{1}{3}\gamma \mu_0 n.
\]
Since the refutation width is $\Width(\Pi) = \max_{C \in \Pi} |C| \ge |C^*|$, and this holds for all valid refutations $\Pi$, the minimal refutation width is lower-bounded by $\alpha n$ where $\alpha = \frac{1}{3}\gamma \mu_0 > 0$.
\end{proof}
